\section{Introduction}\label{sec:introduction}
An algebraic consequence can be simple to state and difficult to derive.
Given polynomial equations $\Gamma_n$ on a universe of size $n$, we ask
how much degree is needed to prove a prescribed equation $G^{(n)}=0$ from
them. Both the axioms and the target are generated by fixed rules as
$n$ varies. The answer has four mutually exclusive and exhaustive
patterns: bounded degree, linear hardness, a periodic mixture of the
two, and systematic fluctuation with finite linear peaks. This paper
proves the classification for uniform polynomial consequences, with
refutation as the special target $G^{(n)}=1$. It strengthens the earlier
logarithmic Nullstellensatz and double-logarithmic Polynomial Calculus
lower bounds to linear bounds in the universe size.

The NS proof system was introduced by Beame, Impagliazzo,
Kraj\'i\v cek, Pitassi and Pudl\'ak~\cite{NSFoundations}; PC was
introduced by Clegg, Edmonds and Impagliazzo~\cite{PCFoundations}.
To fix the proof-theoretic meaning of a consequence, let
$\Gamma=\{f_1,\ldots,f_m\}$ be a finite set of polynomials over a field
$K$, interpreted as the equations $f_i=0$. A \emph{Nullstellensatz}
(NS) derivation of $G$ is an identity
\[
                         G=\sum_{i=1}^m a_i f_i.
\]
Its degree is the largest total degree of a nonzero product $a_i f_i$.
\emph{Polynomial Calculus} (PC) derives polynomials line by line: each
line is an input axiom, a linear combination of earlier lines, or an
earlier line multiplied by a variable. A derivation of $G$ ends with
that actual polynomial, and its degree is the largest total degree
of any line after collecting coefficients. Thus $x-y$ and $y$ derive
$x$ at degree one, while deriving $1$ certifies that the input equations
have no common zero. Soundness makes every derivation a semantic
consequence. The converse is not imposed as an inference rule:
semantic entailment and the existence of an NS or PC derivation are
distinct questions for a general polynomial system.

Degree measures the complexity of the algebraic calculation rather
than its number of lines. A fixed degree bound confines all lines to
a vector space with polynomially many monomials in the variable count.
An NS certificate uses the span of the allowed monomial multiples of
axioms. PC can also cancel terms and multiply the resulting polynomial
again. Its bounded-degree consequences therefore form an iterated
linear closure, and its minimum degree can be smaller than the NS
minimum. Section~\ref{sec:systems} gives the exact characterisations
used throughout the proof.

A \emph{uniform family} $(\Gamma_n)_n$ is a family of systems of
polynomial equations, generated by a fixed finite list of bounded-degree
polynomial patterns using the same coefficient rules at every universe
size. At each size $n$, all equations in $\Gamma_n$ are imposed
simultaneously. Each pattern can generate many equations as its indices
range over the universe, so the number of equations can grow with $n$
even though the list of patterns is fixed.
Variable and axiom indices each carry a bounded number of labels,
and relabelling the universe permutes the equations. Each axiom is fixed
by permutations fixing its distinguished index labels, even when it
contains sums over the entire universe. A \emph{fully invariant uniform
target} has no distinguished labels and is unchanged by every
permutation; $G^{(n)}=\sum_{i<n}x_i$ is a simple example. Its logical
antecedents include the Paris--Wilkie translation~\cite{ParisWilkie},
treated systematically by Kraj\'{i}\v{c}ek~\cite[Chapter~9]{KrajBounded},
and my early work on independence and finitisation in bounded
arithmetic~\cite{RiisThesis,RiisFinitisation}, followed by the
complexity-gap results~\cite{RiisTree,Riis}. My joint journal papers
with Meera Sitharam develop the uniform translation and its
permutation-module viewpoint~\cite{RiisSitharamTautologies,RiisSitharamModules}.
Here the hypotheses are stated entirely in algebraic terms. Definitions~\ref{def:template}
and~\ref{def:uniform-target} make the coefficient rules and support
conditions precise.

\paragraph{Examples of uniform families.}
For these examples label the universe by $1,\ldots,n$. With three
unary variable types $x_i,y_i,z_i$, consider the generating polynomial
\[
 P_{1,n}=x_1+y_1+\sum_{j,k=1}^{n}x_jx_k+
                       \sum_{i,j,k=1}^{n}z_i z_j z_k.
\]
Its total degree is at most three, independent of $n$, although the
number of summands and variables grows with $n$. A permutation fixing
$1$ fixes this whole polynomial. Its distinguished index therefore has
support $\{1\}$. Taking its permutation orbit gives the symmetric axiom
system $P_{a,n}=0$ for all $a\in\{1,\ldots,n\}$, where only the first
two terms become $x_a+y_a$: both global sums are unchanged.
The degree bound concerns each monomial, not the number of terms or
the union of labels mentioned by the sum. All sums are evaluated in
the fixed coefficient field, including any coefficient cancellations.

Binary and higher-arity variable types are equally admissible. For
binary variables $u_{ij}$, the row axioms
\[
               F_{i,n}=\sum_{j=1}^{n}u_{ij}-1=0
\]
have degree one and index support $\{i\}$, while each variable has
support contained in $\{i,j\}$. A cubic polynomial such as
\[
       H_n=\sum_{i,j,k=1}^{n}u_{ij}u_{jk}u_{ki}
\]
is fully invariant, with one empty-support axiom or target index.
A single system can combine several such patterns. For example, it can
impose all $n$ row equations $F_{i,n}=0$, all $n^2$ Boolean equations
$u_{ij}^2-u_{ij}=0$, and the global equation $H_n=0$ simultaneously.
These three patterns define a system of $n+n^2+1$ indexed equations.
There can be $n^2$ binary variables, or $n^a$ variables of any fixed
arity $a$, without violating the bounded-support hypothesis. These
examples are instances of the rank-independent coefficient rules in
Definition~\ref{def:template}, rather than arbitrary symmetric
polynomials chosen afresh at each rank.

\paragraph{Familiar examples from the literature.}
Algebraic encodings of pigeonhole and counting principles provide
important examples of this uniform structure. Razborov studied
Polynomial Calculus degree lower bounds for pigeonhole
principles~\cite{RazborovPC}; Beame and Riis studied Nullstellensatz
bounds for onto-pigeonhole principles~\cite{BeameRiis}. Fixed-block-size
partition principles also underlie the Nullstellensatz lower bounds of
Beame, Impagliazzo, Kraj\'{i}\v{c}ek, Pitassi and
Pudl\'{a}k~\cite{NSFoundations}. In their bounded-degree encodings,
Boolean constraints, pairwise incompatibility constraints and sums
expressing that each element is assigned exactly once are generated by
fixed patterns invariant under relabelling. Fixed-offset pigeonhole
families fit the present single-universe formulation by treating the
extra elements as permanent parameters, as described below. Thus these bounded-degree, fixed-pattern encodings are instances of
the uniformly generated polynomial systems considered here.
The encoding matters: growing-degree axioms are outside the
bounded-degree hypothesis used here.

\paragraph{Orbit generators and permanent parameters.}
A label in an orbit generator is moved when the orbit is taken. Thus
the orbit of $x_1y_2=0$, with $1\ne2$, consists of $x_ay_b=0$ for
all $a\ne b$. It does not include the diagonal equations. This differs
from permanently naming two distinct domain elements: the single
constraint involving those named elements need not impose the analogous
constraint on every pair.

Permanent parameters can be included by using separate variable types.
For one named element $c$, a binary family $z_{ij}$ splits into the
scalar variable $z_{cc}$, the unary families $z_{ci}$ and $z_{ic}$,
and the binary family $z_{ij}$ for $i,j\ne c$. The letter $c$ is part
of the type name; only the other labels are permuted. Correspondingly,
\[
 \sum_{i,j\in\{c\}\cup U}z_{ij}
 =z_{cc}+\sum_{i\in U}z_{ci}+\sum_{i\in U}z_{ic}
                                      +\sum_{i,j\in U}z_{ij}.
\]
The scalar type has empty support and is fixed by every permutation.
Such zero-arity variable types are allowed by
Definition~\ref{def:template}: the positive parameter $k$ is an upper
bound on support size, not a lower bound on arity. More generally,
a fixed set of $\ell$ distinct named elements splits each finite-arity
variable family into finitely many types, according to which coordinates
are named and which remain free. Sums split in the same way. In a
universe of total size $n$, the remaining action is $S_{n-\ell}$;
reindexing the free universe puts the system in the usual $S_m$ form
with $m=n-\ell$. This fixed shift preserves bounded degree and the
asymptotic linear bounds; statements about residues and initial universes
are read with the same shift. For example, with permanent parameters
$c_1,c_2$, the equation $x_{c_1}y_{c_2}=0$ is a single equation in two
scalar variable types. Its orbit under permutations of the remaining
labels contains only itself.
Here the parameters are preselected distinct domain elements. An
uninterpreted constant symbol, whose value is chosen as part of a model,
is instead represented by its function-graph variables in the logical
translation below.

The proof is organised around one question. If an invariant uniform
target $G^{(n)}$ has a degree-at-most-$d$ derivation from $\Gamma_n$, when
can the corresponding target $G^{(n+q)}$ be derived within the same
degree bound from $\Gamma_{n+q}$? Here $q=p^a$ is the number of additional
elements in the finite universe, whose size increases from $n$ to $n+q$. Once such a lift is available and the universe is sufficiently large
relative to the degree bound, it can be repeated:
\[
 \Gamma_{n+\ell q}\vdash_{\mathrm S,\,\le d}G^{(n+\ell q)}
                    \qquad(\ell=0,1,2,\ldots).
\]
The selected numerical degree $d$ is constant on this progression;
the proof lines and their coefficients may change. The essential
quantitative question is how large the initial universe must be relative
to $d$. For indexed targets the same question is asked for the whole
prescribed target family, using its coefficient map rather than a
literal lift of each polynomial.

\paragraph{First-order logic and finite structures.}
The framework includes polynomial encodings of arbitrary fixed
first-order sentences over a finite vocabulary of uninterpreted
relations and functions. A sentence can first be Skolemised and
normalised to universal clauses, with auxiliary relation or function
symbols where needed. The new functions record existential witnesses:
for example, $\forall u\,\exists v\,R(u,v)$ can be represented by
$\forall u\,R(u,f(u))$. This preserves existence of models on each
nonempty finite universe. A model may have several witness expansions,
so it does not assert a one-to-one correspondence with models of the
original sentence before those witnesses are chosen.

Relation variables record truth values, and function-graph variables
record chosen outputs. Boolean equations, totality equations and
pairwise functionality collisions give the translation of
Section~\ref{sec:systems}. Its polynomial system has a common root
at size $n$ if and only if the original sentence has an $n$-element
model; roots correspond to models of the chosen normalised expansion.
The vocabulary, arities, clause lengths and hence polynomial degrees
are fixed while $n$ grows. This is the translation setting used in
LICS~2008~\cite[Section~1.2, p.~273]{Riis} and the expanded
2009 report~\cite[Section~1.2, pp.~3--4]{RiisLong}, following the joint work
with Meera Sitharam~\cite{RiisSitharamTautologies}; a precise translation
statement is Theorem~1 of its BRICS version~\cite[Section~II]{RiisSitharamTranslationPreprint}. The joint work gives the
translation for existential second-order sentences, including these
first-order existence problems, and its permutation-module setting is
developed in~\cite{RiisSitharamModules}.

Existential quantification is not simply replaced by a sum. For a
witness graph, $\sum_j x_{ij}=1$ is used together with Booleanity and
$x_{ij}x_{ij'}=0$ for $j\ne j'$. Only together do these require exactly
one chosen output; in characteristic $p$, the sum equation alone also
allows $p+1$ Boolean ones. The example of an injective function omitting
one named point is obtained in this way and has no finite nonempty
model, although it has an infinite model.

The domain symbols here are uninterpreted: order, arithmetic and
other relations are themselves part of the structure being sought.
One may axiomatise a linear order and then quantify over its possible
interpretations. Fixing an external order or arithmetic on the labels
in advance is a different translation and need not satisfy the full
permutation and rank-independent coefficient hypotheses. The theorem
does not automatically cover such built-in structures.

A fixed coefficient rule can likewise prescribe the polynomial
expressions to be derived. At universe size $n$, write these targets as
$G_b^{(n)}$, where $b$ is the target index. The superscript records the
universe size; it does not name a distinguished element. For a fully
invariant target we omit $b$ and write $G^{(n)}$. An indexed target family
may request $G_b^{(n)}=0$ for every $b$, or the representatives of one
fixed index orbit. No arbitrary choice of a different orbit at each
size is made.

The governing hypotheses are fixed bounds on axiom and target degrees,
bounded supports of the
variable and axiom indices, and coefficients prescribed consistently
across universe sizes and under relabelling. An index supported on a
few labels can specify a polynomial involving the entire universe:
the totality equation $\sum_j x_{ij}-1=0$ has index $i$, even though
the sum uses every $j$. The same distinction applies to targets.
Uniformity does not require all terms of a target polynomial to use
one bounded set of labels. The parameter $n$ is universe size, not
variable count; variables indexed by $k$-tuples may number on the
order of $n^k$. The coefficient field has one fixed positive
characteristic throughout.

There is a proof-size interpretation when each polynomial is represented
by its \emph{full, explicit coefficient list}. For a polynomial of
total degree $d$, this lists all monomials of degree at most $d$ in
the entire ambient variable set, including zero coefficients. The
convention applies to PC lines and NS multipliers as well as to the
input. Bounded-degree derivations can then be written in polynomial
size, whereas linear degree lower bounds force size at least
$2^{\Omega(n)}$ on the relevant instances. These are lower bounds;
no universal matching exponential upper bound is asserted. The size
conclusions concern finite derivations of the prescribed targets.
An underivable target has infinite degree and proof size. Over a fixed
finite field, coefficient counts also give the corresponding bit-size
bounds. If the input length is $L(n)=O(n^{c_0})$ for a fixed $c_0>0$, the exponential bound in
$n$ becomes a stretched-exponential bound
$2^{\Omega(L(n)^{1/c_0})}$. Corollary~\ref{cor:dense-types} supplies the
counting argument and the precise scope; sparse representations that
omit zero coefficients use a different size measure.

\paragraph{Complexity gaps and the earlier classification.}
The restriction to a fixed uniform description belongs to a broader
programme of complexity gap theorems. Resolution combines
$A\lor x$ and $B\lor\neg x$ to derive $A\lor B$. In tree-like
resolution, a derived clause cannot be shared between branches;
general resolution permits this reuse. Riis's tree-resolution gap
theorem~\cite{RiisTree} gives, for the standard translations of a
first-order sentence with no finite models, either polynomial-size
refutations or a $2^{\Omega(n)}$ size lower bound. The hard case occurs
exactly when the sentence has an infinite model; if it has no model
at all, the polynomial case holds. Distinguishing the two cases is
undecidable.

Dantchev and Riis~\cite{DantchevRiis} proved a size gap between
polynomial upper bounds and $2^{\Omega(n^\varepsilon)}$ lower bounds,
for some fixed $\varepsilon>0$, for \emph{full resolution} on
unary-relativised families. A unary predicate selects the domain on which the sentence
is interpreted. When the original sentence has no finite models, its
relativised translations have polynomial-size refutations, even
tree-like ones, if it has no model at all. If it has an infinite model,
every full-resolution refutation has size
$2^{\Omega(n^\varepsilon)}$ for some $\varepsilon>0$. That work also
treats tree-like resolution with bounded-size conjunctions and shows
that a built-in order destroys the gap phenomenon.

Dantchev, Martin and Szeider~\cite{DMS} refine the tree-resolution
classification in the parameterised setting. Dantchev and
Martin~\cite{DMRank} obtain a constant-versus-polynomial rank gap for
the Lov\'asz--Schrijver and Sherali--Adams systems, again governed by
infinite models. Beyersdorff, Clymo, Dantchev and Martin~\cite{QBFGap}
establish tree-like QBF-resolution analogues, with explicit dependence
on the translation. These results concern different resources---size,
parameterised size or rank---and do not themselves yield an NS or PC
degree bound.

Positive characteristic allows more than a dichotomy for algebraic
degree. Easy and hard sizes can occupy different residue classes, and
one family can have infinitely many recurring finite degree levels.
My LICS 2008 paper~\cite[Theorems~1 and~10]{Riis}, with fuller treatment
in the 53-page 2009 report~\cite{RiisLong}, established four refutation types,
their periodic mixed case and agreement of the NS and PC type labels.
Its general degree lower bounds were $\Omega(\log n)$ for NS and
$\Omega(\log\log n)$ for PC. Their periodicity and degree-dependent
onsets came from Kraj\'i\v cek's Theorem~3.5~\cite{Krajicek}, building
on Ajtai's analysis of symmetric linear systems~\cite{Ajtai}.
Here the linear bounds are proved for
the more general task of deriving uniform targets. The constant
target $1$ recovers the refutation classification and its strengthened
bounds.

In the final remarks of the LICS paper and the longer report, I
conjectured that the classification should admit a degree lower bound
of the form $\Omega(n^\varepsilon)$, for sufficiently small fixed $\varepsilon>0$,
at least for NS and possibly for PC~\cite[Section~8]{Riis}\cite[Section~9]{RiisLong}.
I recorded the considerable effort already spent on this question.
Both the LICS abstract and the longer report explicitly asked about a
linear degree lower bound~\cite{Riis,RiisLong}.
The present theorem proves that bound for both systems, with the
hard-residue and finite-peak qualifications below, and extends the
setting from refutations to uniform polynomial consequences.

The journal work on uniformly generated submodules concerns fields of
characteristic zero~\cite{RiisSitharamModules}. In a separate 1997
preprint with Meera Sitharam~\cite{RiisSitharam1997}, we announced a
nonconstant-to-linear degree gap for uniform NS and PC families.
In my 2009 report, I explicitly acknowledged an unresolved gap in the
inducing-down proof and stated that the main claim remained conjectured
\cite[p.~50, footnote~3; p.~53, reference~33]{RiisLong}.
That qualification concerns the proof; it does not refute the asserted
gap theorem. The present argument gives independent finite algebraic
constructions for the hypotheses and four profiles stated here.

The preliminary fresh-label construction---a construction introducing
additional labels---in Section~\ref{sec:quadratic-route} already yields an $\Omega(\sqrt n)$
degree gap in characteristic two, answering the positive-power question
in that setting. The linear construction makes a structural change:
corrections at successive tuple lengths reuse one retained pool of labels.
The odd-characteristic argument supplies the additional nilpotent-chain
structure needed to obtain the same linear degree lower bound in every
positive characteristic.

The linear degree lower bound is asymptotically optimal in general. The
injection example in Section~\ref{sec:examples} has minimum NS and PC
degrees $\Theta(n)$, while the fluctuating example has finite
$\Theta(n)$ peaks. These examples show that a superlinear degree lower
bound cannot replace the linear bound throughout the classification.

\paragraph{The four degree profiles.}
For either proof system, let $h(n)$ be the least derivation degree of
the uniform target; for an indexed target family, take the maximum of
the least degrees over its finite rank-$n$ index set. Infinity records
an underivable target, and an empty index set has degree zero. Let
$\rec(h)$ be the set of finite degrees occurring at infinitely many
sizes. The four alternatives and their coefficient-list size readings
are as follows.
\begin{center}
\begin{tabular}{@{}p{0.065\textwidth}p{0.46\textwidth}p{0.39\textwidth}@{}}
\toprule
Type & Degree behaviour & Derivation-size consequence\\
\midrule
$\mathrm T_1$ & Bounded degree beyond a fixed initial range
 & Polynomial upper bounds.\\
$\mathrm T_2$ & No recurring finite degree; linear lower bounds
 & Exponential lower bounds.\\
$\mathrm T_3$ & Finitely many recurring levels without eventual boundedness;
fixed low degrees on a nonempty proper set of residues, linear lower
bounds on the others
 & Periodic polynomial upper bounds and exponential lower bounds.\\
$\mathrm T_4$ & Infinitely many recurring levels; finite liminf
and finite linear peaks
 & Polynomial subsequences and finite exponential peaks at arbitrarily
large sizes.\\
\bottomrule
\end{tabular}
\end{center}
Finite initial exceptions are allowed. In $\mathrm T_3$, easy residues
need not share the same degree. In $\mathrm T_4$, a linear degree bound
at every sufficiently large size is false: the least recurring finite
degree gives arbitrarily late low-degree derivations. Its infinitely many recurring finite degrees distinguish it from
$\mathrm T_3$, which has only finitely many. Thus arbitrarily large
sizes have finite linear peaks, while one fixed finite degree also
occurs arbitrarily late.
Section~\ref{sec:classification} proves exhaustiveness and exclusivity
and states the exact quantitative bounds.

\paragraph{From unchanged-degree lifting to a linear bound.}
Fix a proof-degree budget $d$. The central statement says that
derivability within this budget is periodic in universe size after
an onset linear in $d$, with a period also linear in $d$. Consequently,
a finite degree $d$ attained beyond its own onset must recur on an
infinite arithmetic progression. A degree which does not recur can
only occur before that onset. Inverting this restriction gives the
linear lower bound in $\mathrm T_2$ and on the hard residues of
$\mathrm T_3$. In $\mathrm T_4$, the bound on the period locates
representatives of unbounded recurring degrees at sizes proportional
to those degrees, producing finite linear peaks. Both the onset and
period are needed.

The relevant budget counts labels. If variable indices have support
at most $k$ and axiom indices at most $s$, a degree-$d$ monomial
paired with an axiom index uses at most $kd+s$ labels. The fully
invariant target uses this budget. For targets indexed by supports
of size at most $t$, the general proof uses the conservative budget
\[
                             R(d)=2kd+s+t.
\]
The extra factor of two comes from detecting the target image outside
the consequence space by a pairing; both factors of that pairing
must be transported. The fixed target degree affects the initial
rank threshold and excludes budgets below the target's degree.
Above that fixed threshold, a sufficient periodic onset is
$n>2R(d)$, and a period is the least power of the characteristic
exceeding $R(d)$. Thus the extension to indexed targets retains
the linear degree lower bounds.

There are three transport obligations. First, the operation must take
the actual axiom and multiplication maps to those at the smaller
rank. Second, it must preserve the sums, images and intersections
which compute all bounded-degree consequences. This is the difficult
PC requirement: a cancellation can create a new polynomial eligible
for multiplication. Third, it must preserve and reflect whether the
prescribed target belongs to the resulting space. Exact deletion and
a common group-algebra projector meet the first two obligations;
transport of the target coefficient map and nonzero-subquotient
detection meet the third. An abstract isomorphism of spaces alone
would supply none of these map compatibilities.

The construction is developed in stages. In characteristic two, a
first lift uses separate fresh labels at successive tuple lengths;
its quadratic label cost already explains a square-root degree gap.
The linear construction reuses one pool. Its projector is constructed
through tuple lengths $b\le R(d)$,
repairing the current length while leaving the shorter lengths
unchanged. The same retained labels serve all repairs, so the
allocation is controlled by $\max_{b\le R(d)}(2b-1)$ rather than
a sum of separate fresh-label allocations. The proof develops the
square-zero construction in characteristic two, normalisation and
the zero-marginal decomposition, the full rook-forcing induction,
and the top-power construction in arbitrary positive characteristic.
For odd $p$, square-zero identities are replaced by full length-$p$
nilpotent chains and an ordered projector, with a separate transpose
argument. These repairs retain the same linear label allocation.
Section~\ref{sec:transfer} states the exact map requirements and
the resulting transfer, and
Sections~\ref{sec:block-algebra}--\ref{sec:odd} establish its algebraic
and combinatorial ingredients.

The comparison between the two systems is separate. Compressing
invariant-subspace chains to a bounded universe gives
$h_{\PC}(n)\le h_{\NS}(n)\le F(h_{\PC}(n))$ whenever the PC
degree is finite, with a monotone function $F$ independent of
universe size. It applies to
the entire consequence space and hence to each prescribed target.
This establishes common type without asserting equality of exact
degrees or of recurring degree sets. The linear PC bound is proved
directly, rather than inferred by inverting the potentially large
function $F$.

\paragraph{Examples, semantics and undecidability.}
The first-order examples illustrate the four refutation types. A
fresh-variable construction also produces a nonconstant target
whose finite derivation degree has linear peaks; this shows that
the consequence theorem includes actual polynomial targets as well
as the constant $1$. The motivating Boolean translation uses
totality sums and functionality collisions, with no automatic
reduction modulo Boolean equations. Section~\ref{sec:model-theory} gives a $p$-adic hardness spectrum for
all uniform requests: the four types correspond to an empty spectrum,
the whole space, a nonempty proper clopen set, and a closed set that is
not open. This makes explicit the residue-class structure already
implicit in my earlier report~\cite[Section~4]{RiisLong}, now in the
setting of uniform consequences. The same section distinguishes
linear witnesses from actual solutions and recalls Kraj\'i\v cek's
Euler criterion. The topological argument needs no Boolean semantic
hypotheses and does not assert a converse to that criterion.

The undecidability theorem already holds in the refutation special
case. The no-finite-model promise and undecidability for every
nonempty proper union of types were stated in LICS~2008
\cite[Theorem~1]{Riis}. Here explicit reductions separate any two
types under that same promise, refining the published every-union
formulation. On normalised first-order inputs, an explicit finite
computation prefix gives a uniform degree bound when its machine
halts, and a controlled disjunction supplies the pairwise reductions.
The effective reduction and its finite-term certificate are included.

\paragraph{Algebraic antecedents.}
Kraj\'i\v cek's bounded-degree NS--PC comparison
\cite[Theorem~5.5]{Krajicek}
supplied an input to my earlier classification. The present proof
retains the NS-span stabilisation argument and supplies its uniform
chain bound independently. The behaviour of fixed-degree data as the
number of labels grows is a central theme of representation stability
and FI-modules, developed by Church, Ellenberg and Farb~\cite{CEF}.
An FI-module assigns vector spaces and compatible linear maps to finite
sets and injections; it organises comparisons as the label set grows. In positive
characteristic, Harman~\cite{Harman} proves periodic equivalences for
bounded-weight symmetric-group representations with a linear stable
range. His equivalences already respect induction, restriction, tensor
products and duality~\cite[Propositions~2.7 and~2.12]{Harman}.
The deletion functor $\ker z/(\ker z\cap\im z)$ used here is also
one of the explicit functors in Hone and coauthors~\cite[Theorem~A]{HKPSYZ}.
Its group-algebra operator comes from a generic shifted cyclic subgroup.
Our argument proves exactness on its stated bounded-support category,
without invoking a categorical equivalence. These antecedents place
the linear-onset mechanism in an established setting. The
present proof independently constructs exact deletion on the actual
polynomial diagrams, preserves their consequence closures and detects
the prescribed targets, with the quantitative bounds used here. The rook
realisation step is related to the integral tensor-space centraliser
results of Bowman, Doty and Martin~\cite{BDM}; the required realisation
is obtained by the forcing and elimination argument in the paper.
Uniform target membership already belongs to Kraj\'i\v cek's
framework~\cite{KrajEuler,Krajicek}. The contributions here are the
linear quantitative bounds for uniform consequences, the independent
algebraic proof and explicit reductions for every pair of types under
the earlier no-finite-model promise.

\paragraph{Notation and proof guide.}\label{sec:specialist-tools}
The principal notation is collected here; local auxiliary symbols are
defined in their respective constructions.
\begin{center}
\begin{tabular}{@{}ll@{}}
\toprule
Notation & Meaning\\
\midrule
$n$, $N=n+q$ & Retained and larger universe sizes\\
$d$; $D$, $D_G$ & Proof-degree budget; fixed axiom and target degree bounds\\
$k,s,t$ & Variable, axiom-index and target-index support bounds\\
$R(d)=\kappa d+\sigma$ & Support budget used in periodicity\\
$N(d)$, $Q(d)$ & Degree-dependent onset and prime-power period\\
$I_d(n)$, $C_d(n)$ & NS and ordinary-PC consequence spaces\\
$P_b(N)$, $\mathcal C_r(N)$ & Injective-tuple modules and their subquotient category\\
$\mathcal Z_b(N)$, $J_r(N)$ & Zero-marginal space and tuple annihilator ideal\\
$z$, $w$, $e$, $\Del$ & Block operator, common inner inverse, projector, deletion\\
\bottomrule
\end{tabular}
\end{center}
The central dependency chain is
\[
 \begin{gathered}
 \text{rook realisation and layer splitting}\Longrightarrow
 \text{common inner inverse}\Longrightarrow\text{exact deletion},\\
 \text{exact deletion and target detection}\Longrightarrow
 \text{periodicity with linear onset}\Longrightarrow\text{four profiles}.
 \end{gathered}
\]
The NS--PC comparison additionally uses a uniform subspace-chain bound.
The matching obstruction supplies the concrete examples needed for
pairwise undecidability. Standard background is in Lang~\cite{LangAlgebra},
Alperin~\cite{AlperinLocal}, the Stacks Project~\cite{StacksExact} and
Cutland~\cite{Cutland}; the specific algebraic and combinatorial
arguments used here are proved below.

\paragraph{Outline of the paper.}
Section~\ref{sec:systems} defines the proof systems, uniform axiom
and target templates, first-order encoding and bounded-degree
consequence spaces. Section~\ref{sec:classification} derives the
four profiles and their size consequences from quantitative transfer.
Sections~\ref{sec:transfer}--\ref{sec:odd} construct the deletion and
lifting operations, first in characteristic two and then in arbitrary
positive characteristic. Section~\ref{sec:comparison} proves uniform
NS--PC type agreement. Section~\ref{sec:model-theory} gives the topological
classification and discusses algebraic witnesses and Euler structures.
Section~\ref{sec:target-examples} develops nonconstant targets and the
fresh-variable construction. Sections~\ref{sec:examples}--\ref{app:matching}
give the four first-order refutation examples and prove their matching
obstruction.
Section~\ref{sec:undecidability} completes the computation-prefix
reduction. Section~\ref{sec:conclusion} closes with the mathematical
consequences and acknowledgements.
